\chapter{项目目标与技术路线}
\section{从指令执行到受约束的物理运动}
机器人接收到目标位置，不代表能够按时、准确且稳定地完成运动。测量采样、传输、驱动响应、接触变化和动力学约束会共同影响执行。本项目从固定基座机械臂出发，先建立可以记录目标与反馈的链路，再研究带外力和支撑切换的运动。

机械臂与双足并非同一个控制问题，但共享几个核心环节：参考生成、状态反馈、受限执行和独立验收。五次轨迹用于连接动作阶段，滤波与预测用于处理测量滞后，接触控制用于协调运动与力，浮动基座控制则必须通过足地反力间接稳定身体。后续扩展是在逐步检验这些接口，而不是将多个演示拼成一个算法。
\section{导师任务与自主扩展}
原月度计划包括仿真平台、IMU状态反馈、机械臂柔顺控制和基于倒立摆的稳定行走分析。正文按技术关系组织，不将章节编号当作周次。倒立摆阶段以月度计划大方向开展，缺失的详细任务书不被推定为已逐项验收。三维人形控制属于在上述基础上的自主扩展。
\begin{longtable}{p{.19\linewidth}p{.35\linewidth}p{.35\linewidth}}
\toprule 计划主题 & 已开展工作 & 报告位置\\\midrule\endhead
平台与轨迹 & ROS2桥接、模型限位、五次轨迹、实测误差与增益对照 & 第2章\\
状态反馈 & 虚拟与公开IMU、实体采集、因果滤波、发布频率与数据年龄 & 第3章\\
柔顺控制 & 法向导纳、接触冲击、连续逆解、重力和手部碰撞、执行接口 & 第4、5章\\
稳定行走分析 & 倒立摆、受限落脚、频率、噪声、推扰与延迟预测 & 第6章\\
自主扩展 & 平面多关节连续行走、统一接触接口、G1三维全身控制 & 第7、8章\\
\bottomrule\caption{任务主题与技术篇章的对应关系}\end{longtable}

针对导师关于重力、手部碰撞、恒力窗口和网络接口的建议，后续保留了无补偿失败、刚性与柔顺材质、物理步长及真实本机UDP对照。针对早期倒立摆和MuJoCo站立验证不连通的问题，后续建立实际摆腿和关节力矩闭环。正文分别保留早期结论与后期修正，不把后来的结果写回旧实验。
\section{模型层次及其边界}
\begin{longtable}{p{.18\linewidth}p{.24\linewidth}p{.24\linewidth}p{.23\linewidth}}
\toprule 层次 & 状态与执行 & 物理条件 & 能够说明的问题\\\midrule\endhead
Panda轨迹 & ROS2关节反馈；位置／速度目标 & Isaac物理驱动，固定基座 & 指令链路与轨迹响应\\
Panda接触 & 末端力与关节状态；力矩 & 探针或手部接触；分组重力设置 & 法向柔顺与冲击\\
线性倒立摆 & 质心位置速度；地面支撑点 & 恒高质点、有限支撑区域 & 模型内平衡与落脚\\
有限脚掌站立 & 刚体姿态、真实地面接触 & 踝关节与有限足底 & 站立受扰恢复\\
平面双足 & 9速度自由度，六电机力矩 & 基座三自由度；实际腿惯量 & 矢状面连续步行\\
三维G1 & 35速度自由度，29电机力矩 & 完整六自由度浮动基座 & 横向平衡与准静态步行\\
\bottomrule\caption{从简化模型到三维动力学的验证层次}\end{longtable}

简化模型的落脚位置可以直接切换，物理腿却需要有限时间和加速度。规划器中的负载比例是可行分配，实际脚地反力由模拟器求解。二者只有在同一动力学系统中建立接口并核对反馈后，才构成完整物理闭环。
\section{测量、时间与统计口径}
\subsection{四种频率不能混用}
物理积分频率、控制更新频率、传感器产生新样本的频率、消息发布频率分别记录。25赫兹传感器可以被200赫兹发布器反复读取，但不会产生八倍的新信息。仿真时间满足100赫兹也不等于主机运行满足100赫兹；渲染负载与进程调度必须另外测量。
\subsection{原始信号与处理信号分别计量}
接触力区分1毫秒原始值、固定10毫秒分箱均值、控制器滤波值和目标值。平滑后的低误差不能代替原始冲击。轨迹误差以同一状态时刻的参考减实际值计算，拟合滞后另外给出，不通过移动参考来降低原始误差。
\subsection{条件数不等于部署成功概率}
本报告的多数结果来自预设参数网格，包含重复配置、确定性重跑和主动制造的失败。相同配置跨组复用时去重，视频重录不增加样本数。不同模型、周期和目标下的通过率不合并为总体成功率；失败轨迹只统计真实终止前的数据，不与完整轨迹直接比较平均误差。
\section{环境与开源实现的使用}
Isaac Sim独立环境与系统ROS2分别运行；MuJoCo2.3.7用于早期接触和平面模型，三维G1使用独立MuJoCo3.2.3环境。每个阶段的具体版本、依赖和命令在正文及附录中给出。

机器人描述来自官方Panda资源或MuJoCo Menagerie；刚体动力学由模拟器计算，二次规划采用OSQP或Clarabel。本项目实现控制目标、接触约束、时序、实验与核验，不重新实现底层物理引擎或数值优化器。TSID、PlaCo和公开LIPM控制器用于设计参考，具体采用范围在相应章节说明。
\section{与学习策略及真实部署的关系}
动作块窗口、时间戳、限幅和超时测试面向未来学习策略的执行接口，但当时动作由解析轨迹产生，没有运行视觉语言动作模型。G1控制器也不是下载的学习步态策略，而是基于动力学与任务约束的控制实现。

真实硬件记录限于第3章明确标注的IMU与Jetson。其余控制使用仿真状态或人工生成的延迟噪声。当前没有企业设备代测、完整状态估计、真实电机带宽或硬实时验证，报告不以仿真结果替代这些项目。


% BEGIN SIM_VISUAL_OVERVIEW
\section{主要仿真场景与阅读顺序}
以下三幅图先说明实验对象。后续各章同时给出相机画面和定量记录；模型示意、物理录屏和已保存状态重建分别标注。

\motionfigure{figures/ch05/panda_gravity_comparison_full_frame.png}{Panda固定基座接触：比较夹爪与平面的实际相对运动。}{fig:overview-panda}

\motionfigure{figures/ch07/planar_twenty_steps_full_frame.png}{平面双足连续行走：躯干、双腿和实际脚地接触共同推进。}{fig:overview-biped}

\motionfigure{figures/ch08/g1_walk_dual_camera_180.png}{G1三维浮动基座：斜前方和正前方视角展示横向重心转移。}{fig:overview-g1}

% END SIM_VISUAL_OVERVIEW
